\documentclass[10pt]{article}
\usepackage[margin=0.8in]{geometry}
\usepackage[T1]{fontenc}
\usepackage{lmodern}
\usepackage{amsmath,amssymb,amsthm}
\usepackage{graphicx}
\usepackage[hidelinks]{hyperref}

\newtheorem{theorem}{Theorem}
\newcommand{\R}{\mathbb R}

\title{A Two-Point Counterexample to a\\
Nonlinear Lipschitz-Summing Factorization Converse}
\author{Automated mathematical discovery run \texttt{fa\_banach\_001}\\
Lane 6, \texttt{agent\_lane\_06}}
\date{9 August 2026}

\begin{document}
\maketitle

\noindent\textbf{Status.}
Candidate counterexample, likely valid, subject to human review.

\section{Source problem}

Botelho, Maia, Pellegrino, and Santos consider an arbitrary nonempty set
$X$, a metric space $Y$, a compact Hausdorff space $K$, an arbitrary map
$\Psi:X\to C(K)$, and a map $u:X\to Y$ \cite{source}. They call $u$
\emph{$\Psi$-Lipschitz $p$-summing} if some $C\geq0$ satisfies
\begin{equation}\label{eq:summing}
 \left(\sum_{j=1}^m d_Y(u(x_j),u(q_j))^p\right)^{1/p}
 \leq C\sup_{\varphi\in K}
 \left(\sum_{j=1}^m
 |\Psi(x_j)(\varphi)-\Psi(q_j)(\varphi)|^p\right)^{1/p}
\end{equation}
for every finite family $x_j,q_j\in X$.

Corollary~2.5 says that $u$ is $\Psi$-Lipschitz $p$-summing if there are a
probability measure $\mu$ on $K$ and a Lipschitz map
$\widetilde u:L_p(K,\mu)\to Y$ satisfying
\[
 \widetilde u\circ j_p\circ\Psi=u.
\]
On PDF page~6, the authors ask whether the converse holds for $p=2$.

\begin{center}
\includegraphics[width=0.86\textwidth]{figures/open_problem_crop.png}
\end{center}

\section{Proof intuition}

When $K$ is a singleton, the proposed factorization space is forced to be the
connected line $L_p(K)=\R$. The definition of $\Psi$-Lipschitz summability,
however, permits a disconnected target. Map the two points $-1$ and $1$
identically into the two-point metric space. The summing inequality then holds
with equality, but any factorization would continuously map the connected line
onto a set containing both target points. That is impossible.

\section{Counterexample}

\begin{theorem}\label{thm:main}
For every $1\leq p<\infty$, the converse of Corollary~2.5 in
\cite{source} is false. This remains true with $K$ a singleton and $Y$ a
complete two-point metric space. In particular, the source open problem for
$p=2$ has a negative answer.
\end{theorem}

\begin{proof}
Work over the real scalar field. Let $K=\{\ast\}$. Identify $C(K)$ with
$\R$ by evaluation at $\ast$. Put
\[
 X=Y=\{-1,1\},\qquad d_Y(s,t)=|s-t|,
\]
and define
\[
 \Psi(x)(\ast)=x,\qquad u(x)=x.
\]
The target $Y$ is finite and therefore complete.

For arbitrary $m\geq1$ and $x_j,q_j\in X$, we have
\begin{align*}
 \left(\sum_{j=1}^m d_Y(u(x_j),u(q_j))^p\right)^{1/p}
 &=\left(\sum_{j=1}^m|x_j-q_j|^p\right)^{1/p}\\
 &=\sup_{\varphi\in K}
 \left(\sum_{j=1}^m
 |\Psi(x_j)(\varphi)-\Psi(q_j)(\varphi)|^p\right)^{1/p}.
\end{align*}
Thus $u$ is $\Psi$-Lipschitz $p$-summing with constant one.

There is exactly one regular Borel probability measure $\mu$ on $K$, namely
$\mu(\{\ast\})=1$. Evaluation at $\ast$ identifies $L_p(K,\mu)$ isometrically
with $\R$, and under this identification $j_p\Psi(-1)=-1$ and
$j_p\Psi(1)=1$.

If the converse of Corollary~2.5 held, there would be a Lipschitz map
$\widetilde u:\R\to Y$ satisfying
\[
 \widetilde u(-1)=-1,\qquad \widetilde u(1)=1.
\]
Every Lipschitz map is continuous. Since $\R$ is connected, its image under
$\widetilde u$ must be connected. But the only connected subsets of the
discrete two-point space $Y$ are singletons, whereas the displayed conditions
force the image to contain both points. This contradiction proves the result.
\end{proof}

\section{Scope, verification, and novelty}

The counterexample uses precisely the source hypotheses: $X$ may be an
arbitrary nonempty set and $Y$ an arbitrary metric space. It therefore answers
the printed problem, and actually works for every finite $p$. It does not
address a strengthened Banach-valued or absolute-retract-valued version, for
which the connectedness obstruction is unavailable.

The verification report rechecks the summing identity, uniqueness of the
probability measure, and the topology of the forced extension. A bounded
search on 9 August 2026 covered the run's cheap indexes and parsed arXiv
corpus, plus exact-phrase and close web searches for the source problem. The
search found the source paper and later papers citing it, but no exact answer
or matching counterexample. Novelty confidence is moderate pending expert
review.

\begin{thebibliography}{9}
\bibitem{source}
G. Botelho, M. Maia, D. Pellegrino, and J. Santos,
\emph{A unified factorization theorem for Lipschitz summing operators},
Quarterly Journal of Mathematics \textbf{70} (2019), 1521--1533;
arXiv:1902.02569.
\end{thebibliography}

\end{document}
